\documentclass[a4paper]{ltjsarticle}
\usepackage{amsmath, amssymb, amsthm}

\title{計算数学特論 第2回レポート}
\author{学籍番号：62216580\\氏名：平田智也}
\date{提出日：2025/07/23}

\begin{document}
\maketitle

今回の解答に当たっては、問題理解、並びに私の解答をLaTexの形式に成型する過程で
生成AIであるchatGPTを用いた。理論的な導出部分はなるべく自分の力で取り組んだ
つもりであるが、一部の定理の結果は参考文献による。
また、選択した問題は「A1, B1, C1, C2」とする。当該部分にも記載したが、
C3についても、取り組んだ痕跡は記載した。しかし、不完全であるため、
採点いただかなくても構わない。

\section*{A1}

\subsection*{定義}
$n$ 次正方行列 $A = (a_{ij}) \in \mathbb{R}^{n \times n}$ が
\textbf{二重確率行列} であるとは、
以下の3条件を満たすことである：
\begin{align}
  & a_{ij} \ge 0 \quad (i,j = 1,2,\dots,n), \\
  & \sum_{j=1}^n a_{ij} = 1 \quad (i = 1,2,\dots,n), \\
  & \sum_{i=1}^n a_{ij} = 1 \quad (j = 1,2,\dots,n).
\end{align}

また、\textbf{置換行列} とは、
各行と各列にちょうど1つだけ $1$ をもち、それ以外の成分がすべて $0$ である
$n \times n$ 行列をいう。

\medskip
任意の二重確率行列 $A$ に対して
有限個の置換行列 $P_1,\dots,P_k$ と非負係数 $\lambda_1,\dots,\lambda_k$ が存在し、
\[
  A = \sum_{r=1}^k \lambda_r P_r, \quad \lambda_r \ge 0, \quad \sum_{r=1}^k \lambda_r = 1
\]
と表せることを示す。

\subsection*{証明}
$n$ 次二重確率行列全体の集合
   \[
     B_n = \left\{
       A \in \mathbb{R}^{n \times n} \,\middle|\,
       a_{ij} \ge 0,\,
       \sum_{j=1}^n a_{ij} = 1,\,
       \sum_{i=1}^n a_{ij} = 1
     \right\}
   \]
   は、非負制約と $2n$ 本の線形等式で定義される有界な凸多面体である。

任意の $A B \in B_n$ と $0 \le \theta \le 1$ に対し
\[
	C := \theta A + (1-\theta) B
\]
とすれば、行和・列和ともに、
\[
\theta \times 1 + (1-\theta) \times 1 = 1
\]
より、$C$ も行・列和が 1 で、各成分は非負であるため、$C \in B_n$。

置換行列は定義より明らかに $B_n$ に含まれる。

\textbf{Birkhoff–von Neumann 定理}（\cite{micro0616}）より、
   $B_n$ の極点はすべて置換行列である。

有界な凸多面体は、有限個の極点の凸結合で表現できる。
   よって、任意の $A \in B_n$ は置換行列 $P_r$ を用いて
   \[
     A = \sum_{r=1}^k \lambda_r P_r, \quad \lambda_r \ge 0, \ \sum_{r=1}^k \lambda_r = 1
   \]
   と表すことができる。

以上により、任意の二重確率行列は置換行列の凸結合として表せることが示された。

\section*{B1 (a)}

$n$ 個のタスク集合を $S = \{1, 2, \dots, n\}$ とする。
タスク $i \in S$ は開始時刻 $t_i \in \mathbb{Z}$、終了時刻 $u_i \in \mathbb{Z}$ で与えられ、
このタスクを実行すると利益 $p_i \ge 0$ が得られるとする。
各時刻において同時に実行できるタスクは高々1つであり、
時刻 $u_i$ に終了するタスクと時刻 $t_j$ に開始するタスクは同時に選択可能である。

\medskip
ここで、タスク $i$ を選ぶとき $x_i = 1$、選ばないとき $x_i = 0$ となる
二値変数 $x_i \in \{0,1\}$ を導入する。
目的は、選択したタスクの利益和
\[
  \sum_{i \in S} p_i x_i
\]
を最大化することである。

\subsection*{定式化}
各整数時刻 $t$ に対して、時刻 $t$ に実行されているタスクの数は
\[
  \sum_{i:\, t_i \le t < u_i} x_i \le 1
\]
を満たさなければならない。
これにより、任意の時刻に2つ以上のタスクが同時に実行されることを防ぐ。

したがって、整数計画問題は次のように表せる：
\[
  \begin{array}{ll}
    \text{max}   & \displaystyle \sum_{i \in S} p_i x_i \\[6pt]
    \text{s.t.}  & \displaystyle \sum_{i :\, t_i \le t < u_i} x_i \le 1 \quad (\forall t), \\[6pt]
                 & x_i \in \{0,1\} \quad (i \in S).
  \end{array}
\]

\subsection*{行列 $A$ とベクトル $b$ を用いた表現}
これを行列表現で書くと、行列 $A \in \{0,1\}^{m \times n}$ を
\[
  A_{t,i} =
  \begin{cases}
    1 & \text{if } t_i \le t < u_i, \\
    0 & \text{otherwise},
  \end{cases}
\]
と定義することで
\[
  A x \le b
\]
と書ける。ここで $b \in \mathbb{R}^m$ はすべての成分が $1$ である列ベクトルであり、$m$ は可能な整数時刻の数である。

したがって、この問題は
\[
  \text{max} \{ p^T x \mid A x \le b, \, x \in \{0,1\}^n \}
\]
と表現できる。

\section*{B1 (b)}

次に、行列 $A$ が \textbf{完全単模 (以下 TUM )} であることを示す。

\medskip
まず、TUM の定義を述べる：
整数行列 $A \in \mathbb{Z}^{m \times n}$ が TUM であるとは、
その任意の正方部分行列 $A'$ が
\[
  \det(A') \in \{-1, 0, 1\}
\]
を満たすことである。

\medskip
ここで、行列 $A$ は \textbf{区間行列} になっている。
すなわち、各列 $i$ において、$A_{t,i} = 1$ となる行 $t$ の集合は
\[
  \{ t_i, t_i + 1, \dots, u_i - 1 \}
\]
という連続した区間である。

\medskip
\textbf{補題：} 区間行列は完全単模である。

\medskip
\textbf{証明（Sketch）：}
区間行列 $A$ は、各列における $1$ の出現位置が連続している。
この性質により、Ghouila-Houri の判定基準を満たすことが知られており、
区間行列は完全単模（TUM）である（\cite{unimodularPDF} の Corollary 4 を参照）。

\medskip
以上より、B1 の行列 $A$ は完全単模である。
\qed

\section*{B1 (c)}

(a) で定式化した整数計画問題は、(b) で示したように行列 $A$ が完全単模 (TUM) であるため、
その線形緩和問題
\[
  \text{max} \{ p^T x \mid A x \le b, \, 0 \le x \le 1 \}
\]
の最適解は必ず整数解となる。したがって、整数制約 $x_i \in \{0,1\}$ を外した LP を
標準的な多項式時間アルゴリズムで解くことで、効率的に最適解を得ることができる。

さらに、この問題は区間スケジューリング問題と同型であり、
タスクを終了時刻の昇順に並べ、各タスク $i$ について
\[
  DP(i) = \max \{ DP(i-1),\ p_i + DP(\text{非重複なタスクのうち $i$ の直前}) \}
\]
と定めることで、$O(n \log n)$ 時間の動的計画法によっても効率的に解くことが可能である。

\section*{C1 (a)}

$A \in \mathbb{F}^{m \times n}$ を $\mathrm{rank}(A) = m$ とする。
行基本変形と列置換により
\[
  A = ( I_m \mid A')
\]
と書けると仮定する。

線形マトロイド $M(A)$ は、$A$ の列ベクトル $a_1,\dots,a_n$ により定義され、
$E = \{1,2,\dots,n\}$ 上で
\[
  \mathcal{I} = \{ I \subseteq E \mid \{a_i \mid i \in I\} \text{ が線形独立} \}
\]
と表される。

\medskip
マトロイドの双対 $M(A)^*$ は、基底 $B$ に対し $E \setminus B$ を基底とするマトロイドである。
$A$ が $(I_m \mid A')$ の形なら、最初の $m$ 列が基底であるため、
双対基底は残りの $n-m$ 列で張られる。

ここで
\[
  A^* = ( A'^T \mid I_{n-m} )
\]
を考える。このとき、
\[
  A A^{*T} = I_m A'^T + A' I_{n-m}^T = 0
\]
が成り立ち、$A^*$ の行空間は $A$ の行空間の直交補空間を張る。
したがって、$M(A)^*$ は $A^*$ を行列表現とする線形マトロイドである。
\qed

\section*{C1 (b)}

$A_1, A_2 \in \mathbb{F}^{m \times n}$ をランク $m$ の行列とし、
それぞれの列ベクトルによる線形マトロイドを $M(A_1), M(A_2)$ とする。
$M(A_1)$ と $M(A_2)$ の共通独立集合の最大サイズが $m$ であるとは、
ある $m$ 本の列 $J = \{ j_1, \dots, j_m \} \subseteq \{1,2,\dots,n\}$ を選んで
\[
  \det A_1[J] \neq 0, \quad \det A_2[J] \neq 0
\]
となることと同値である。

ここで、不定元 $x_1, x_2, \dots, x_n$ を対角成分に持つ対角行列
\[
  D = \mathrm{diag}(x_1, x_2, \dots, x_n)
\]
を考える。

行列
\[
  M(x) = \begin{pmatrix}
    A_1 \\
    A_2
  \end{pmatrix} D
\]
の行列式 $\det M(x)$ を $x_i$ の多項式として展開すると、
各 $m$ 列の選択 $J = \{j_1, \dots, j_m\}$ に対して
\[
  \det M(x) = \sum_{J: |J| = m}
    \det A_1[J] \cdot \det A_2[J] \cdot x_{j_1} \dots x_{j_m}
\]
となる。

もし共通独立集合がサイズ $m$ で存在すれば、ある $J$ に対して
$\det A_1[J] \cdot \det A_2[J] \neq 0$ となるため、
$\det M(x)$ は非ゼロの多項式となり、適切な $x_i$ を選べば $M(x)$ は正則である。

逆に、$\det M(x)$ が非零多項式であれば、
ある $J$ で $\det A_1[J] \cdot \det A_2[J] \neq 0$ となるため、
その $J$ は共通独立集合であり、サイズは $m$ となる。

したがって、共通独立集合の最大サイズが $m$ であることと、
行列 $M(x)$ が正則となることは同値である。
\qed

\section*{C2 (a)}

二部グラフ $G = (U, V; E)$ を考える。
ここで、$U = \{u_1, u_2, \dots, u_m\}$、$V = \{v_1, v_2, \dots, v_n\}$ とし、
辺集合 $E \subseteq U \times V$ が与えられている。

$U$ の部分集合 $X \subseteq U$ に対して、次の集合族を定義する：
\[
  \mathcal{I} = \{ X \subseteq U \mid G[X, Y] \text{ はある } Y \subseteq V \text{ に対して完全マッチングをもつ} \},
\]
ここで $G[X, Y]$ は $X$ と $Y$ が誘導する部分グラフである。

これにより、$M = (U, \mathcal{I})$ がマトロイドであることを示す。

(I1)
$X = \emptyset$ のとき、$G[\emptyset, \emptyset]$ は自明に完全マッチングをもつ（辺数0のマッチング）。
したがって、$\emptyset \in \mathcal{I}$。

(I2)
$X \in \mathcal{I}$ のとき、$X$ の各頂点は $V$ の異なる頂点とマッチしている。
任意の $Y \subseteq X$ を取っても、$Y$ の頂点は $X$ のマッチングの部分集合を使ってマッチさせることができる。
したがって $Y \in \mathcal{I}$。

(I3)
$X, Y \in \mathcal{I}$ で $|X| < |Y|$ とする。
$Y$ には完全マッチング $M_Y$ が存在する。
$X$ のマッチング $M_X$ と $M_Y$ を比較すると、$Y \setminus X$ の中には $X$ とマッチングの競合を起こさない頂点が少なくとも1つ存在する。
Hallの定理に基づき、$X \cup \{u\}$（$u \in Y \setminus X$）に対しても完全マッチングを拡張できる。
したがって $X \cup \{u\} \in \mathcal{I}$ が成立する。

以上より、$(U, \mathcal{I})$ はマトロイドである。

\qed

\section*{C2 (b)}

ここでは、ジョブスケジューリング問題をマトロイド構造に基づいて定式化し、最小コスト解を求める。

\subsection*{(i) マトロイドとしての定式化}

$n$ 個のジョブ $J = \{J_1, J_2, \dots, J_n\}$ を1台のマシンで処理する。
各ジョブ $J_i$ の処理時間は1時間、納期は $d_i \in \mathbb{Z}_{>0}$、未処理の場合のコストは $c_i \ge 0$ である。

納期までにジョブ $J_i$ を処理できるとは、
「時間スロット $1, 2, \dots, d_i$ のいずれかに $J_i$ を割り当て可能である」ということを意味する。

この制約を表現するため、次の二部グラフ $G = (U, V; E)$ を構成する：
\begin{itemize}
  \item $U = J = \{J_1, J_2, \dots, J_n\}$ （ジョブ集合）
  \item $V = \{1, 2, \dots, T\}$ （$T = \max_i d_i$ とし、各時間スロットを表す）
  \item 辺 $(J_i, t)$ を $1 \le t \le d_i$ のときに張る。
\end{itemize}

ここで、部分集合 $X \subseteq U$ が「独立」であるとは、
$X$ を $V$ の $|X|$ 個の異なるスロットにマッチングできることと定義する：
\[
  \mathcal{I} = \{ X \subseteq U \mid G[X, Y] \text{ が完全マッチングをもつ} \}.
\]
これは C2(a) で述べたマッチングマトロイドと一致し、$(U, \mathcal{I})$ はマトロイドとなる。

本問題の目的は「総コストを最小化すること」だが、未処理ジョブのコストを最小にすることは、
処理するジョブのコスト総和を最大にすることと同値である。
したがって、次の最適化問題を考える：
\[
  \max \left\{ \sum_{J_i \in X} c_i \ \middle| \ X \in \mathcal{I} \right\},
\]
これは **重み付きマトロイド基問題**であり、貪欲アルゴリズムで効率的に解くことができる。

\subsection*{(ii) 貪欲アルゴリズムと最適解}

表1に与えられたジョブのデータは次の通りである：
\[
\begin{array}{c|c|c}
\text{ジョブ} & d_i & c_i \\ \hline
J_1 & 1 & 10 \\
J_2 & 2 & 9 \\
J_3 & 3 & 7 \\
J_4 & 2 & 6 \\
J_5 & 3 & 4 \\
J_6 & 6 & 2 \\
\end{array}
\]

貪欲アルゴリズムは次のように動作する：
\begin{enumerate}
  \item コスト $c_i$ の大きい順にジョブを並べる。
  \item 各ジョブ $J_i$ について、納期 $d_i$ 以下の空きスロットがあれば、最も遅いスロットに割り当てる。
  \item 割り当て不可の場合は未処理とする。
\end{enumerate}

実際の処理過程は以下の通り：
\[
\begin{array}{c|c|c|c}
\text{順序} & \text{ジョブ} & c_i & \text{割当スロット} \\ \hline
1 & J_1 & 10 & 1 \ (\text{成功}) \\
2 & J_2 & 9  & 2 \ (\text{成功}) \\
3 & J_3 & 7  & 3 \ (\text{成功}) \\
4 & J_4 & 6  & - \ (\text{不可}) \\
5 & J_5 & 4  & - \ (\text{不可}) \\
6 & J_6 & 2  & 6 \ (\text{成功}) \\
\end{array}
\]

したがって、最終スケジュールは
\[
\text{スロット} = [J_1, J_2, J_3, -, -, J_6]
\]
となり、処理するジョブ集合は $\{J_1, J_2, J_3, J_6\}$ である。

未処理ジョブ $J_4, J_5$ に対するコストは
\[
c_4 + c_5 = 6 + 4 = 10
\]
であり、これが目的関数値の最小値となる。

\qed

\section*{C3 (a)}

以下、C3を解答することを試みたが、I3が成り立つかどうかを調べるところで詰まってしまったため、
完全なものになっていない。しかし、一応記載は消さずに残す。

無向グラフ $G=(V,E)$ に対して、
\[
F_1 = \{ F \subseteq E \mid (V,F) \text{ の各連結成分は高々1つのサイクルを含む} \}
\]
と定義する。このとき、$F_1$ がマトロイドの独立集合族であることを示す。

\subsection*{解答}
マトロイドの独立集合族 $\mathcal{I}$ は以下の3条件を満たす必要がある：
\begin{enumerate}
  \item[(I1)] $\emptyset \in \mathcal{I}$。
  \item[(I2)] $I \in \mathcal{I}$ かつ $J \subseteq I$ ならば $J \in \mathcal{I}$。
  \item[(I3)] $I,J \in \mathcal{I}$ で $|I| < |J|$ のとき、ある $e \in J \setminus I$ が存在して $I \cup \{e\} \in \mathcal{I}$。
\end{enumerate}

まず、$\emptyset$ は各連結成分にサイクルを含まないので、(I1) を満たす。

次に、$I \in F_1$ で $J \subseteq I$ をとると、$J$ は $I$ よりも辺を少なく持ち、
$I$ の持つサイクルを $J$ も持ち、それゆえ $J$ はそれ以外のサイクルを持たず、
$J \in F_1$ となる。よって (I2) が成立する。

最後に (I3) を確認する。$I, J \in F_1$ で $|I| < |J|$ とする。
グラフ $(V,J)$ では各連結成分に高々1つのサイクルがあるため、
$|J| - |V| + c(J) \le 1$ である（ここで $c(J)$ は $(V,J)$ の連結成分数を表す）。

仮にすべての $e \in J \setminus I$ について
$I \cup \{e\} \notin F_1$ だとすると、$I$ の各連結成分に既にサイクルが存在し、
さらに $|I| \ge |J|$ となるはずであり、これは $|I| < |J|$ に矛盾する。
したがって、$J \setminus I$ の中に $I$ に加えても
2つ目のサイクルを作らない辺 $e$ が存在し、$I \cup \{e\} \in F_1$ が成り立つ。

このことから、(I3) が成立する。


したがって、$F_1$ はマトロイドの独立集合族である。
さらに、ランク関数 $r_{F_1}$ は各連結成分に対し「頂点数－1（木）」または「頂点数（サイクルを含む場合）」となるため、
\[
  r_{F_1}(X) = |V| - c(X) + s(X)
\]
と書ける。ここで $c(X)$ は $(V,X)$ の連結成分数、$s(X)$ はサイクルを含む連結成分の数（各成分につき高々1）である。
\qed

\section*{C3 (b)}

次に
\[
F_2 = \{ F \subseteq E \mid (V,F) \text{ の各連結成分は高々1つのサイクルを含むが、偶数長サイクルを含まない} \}
\]
がマトロイドの独立集合族であるかを調べる。

\subsection*{解答}
$I \in F_2$ であるとは、$(V,I)$ の各連結成分が「木」または「奇数長サイクルを1つだけ含むグラフ」であることを意味する。

(I1)
$\emptyset \in F_2$ は明らかである。

(I2)
$I \in F_2$ で $J \subseteq I$ を取ると、$J$ の連結成分は $I$ の連結成分の部分グラフであるため、
新たに偶数長サイクルが生じることはない。したがって $J \in F_2$ である。

(I3)

\qed


\section*{C3 (c)}

次に
\[
F_3 = \{ F \subseteq E \mid (V,F) \text{ の各連結成分は高々1つのサイクルを含むが、奇数長サイクルを含まない} \}
\]
がマトロイドの独立集合族であるかを調べる。

\subsection*{解答}
$I \in F_3$ とは、$(V,I)$ の各連結成分が「木」または「偶数長サイクルを1つだけ含むグラフ」であることを意味する。

(I1)
$\emptyset \in F_3$ は明らかである。

(I2)
$I \in F_3$ で $J \subseteq I$ を取ると、$J$ の連結成分は $I$ の連結成分の部分グラフであり、
新たに奇数長サイクルが生成されることはない。したがって $J \in F_3$ である。

(I3)


\qed

\begin{thebibliography}{99}
\bibitem{unimodularPDF}
K. Truemper,
``Unimodular Matrices,'' Lecture Notes, 2010.
\bibitem{micro0616}
A. Abdulkadiroğlu and T. Sönmez,
``Implementing Random Assignments: A Generalization of the Birkhoff–von Neumann Theorem,''
\emph{Econometrica}, vol.~70, no.~2, pp.~687–702, 2002.
\end{thebibliography}

\end{document}
